% Capítulo 1 – gerado a partir do relatório; edite livremente
\chapter{INTRODUÇÃO}

A peça escolhida para este projeto é a engrenagem cilíndrica de dentes retos do conjunto engrenagem e cremalheira apresentado no material do PROTEC (Figura~\ref{fig:desenho}). Trata-se de uma engrenagem de módulo 6 mm, com 24 dentes e ângulo de pressão de 20°, em aço SAE 8620. Engrenada à cremalheira, ela transforma a rotação do eixo em movimento de translação, solução usada, por exemplo, em portões industriais deslizantes, carros de transporte, mesas de máquinas e mecanismos de elevação, nos quais a transmissão trabalha com carga elevada e baixa velocidade.


Para que o engrenamento ocorra sem ruído, com folga controlada e com o contato distribuído ao longo da largura do dente, o dentado precisa ser concêntrico ao furo de montagem, e as faces do cubo, que fazem o encosto axial da engrenagem no eixo, precisam ser perpendiculares a esse furo. Por isso, ao longo de todo o planejamento, o furo e uma das faces do cubo são gerados na mesma fixação e usados como referência até a retificação final.


O material também orienta a sequência de fabricação. O SAE 8620 é um aço de cementação: quase toda a usinagem é feita com o aço recozido e, depois da cementação e da têmpera, apenas as superfícies funcionais são retificadas.


\section{Objetivos}

O objetivo deste trabalho é apresentar a seleção de operações, ferramentas, parâmetros de corte e máquinas-ferramenta para a fabricação da engrenagem a partir de uma barra laminada de aço SAE 8620 com diâmetro de 160 mm, incluindo a retificação, que é obrigatória no projeto. Para cada operação são definidos a fixação da peça e a trajetória das ferramentas, calculadas as potências de corte e estimados os tempos de usinagem, e o resultado é reunido na folha de processo.

